%==============================================================================
\section{Agglomerative Clustering}
\label{sec:agglomerative}
%==============================================================================

\begin{dinhnghia}[title={Agglomerative Clustering}]
Thuật toán phân cụm \textbf{phân cấp bottom-up}: mỗi điểm ban đầu là một cụm, lặp gộp hai cụm gần nhất đến khi đạt tiêu chí dừng.
Sinh \textbf{dendrogram} thể hiện quan hệ phân cấp giữa các cụm.
\end{dinhnghia}

\subsection{Các bước thuật toán}

\begin{phuongphap}[title={Các bước}]
\begin{enumerate}
  \item Coi \textbf{mỗi điểm} là một cụm riêng.
  \item Tính \textbf{ma trận gần} (proximity) bằng metric khoảng cách.
  \item \textbf{Gộp} hai cụm theo tiêu chí linkage (single, complete, average, ward, \ldots).
  \item \textbf{Cập nhật} ma trận khoảng cách.
  \item Lặp bước 3--4 cho đến khi còn \textbf{một cụm} (hoặc dừng sớm theo ngưỡng).
\end{enumerate}
\end{phuongphap}

\begin{ynghia}[title={Vì sao dùng Agglomerative?}]
Dễ diễn giải quan hệ giữa điểm; không bắt buộc biết số cụm trước khi xây dendrogram (số cụm chọn khi ``cắt'' cây).
Hiệu quả, phát hiện được cụm nhỏ.
\end{ynghia}

\subsection{Triển khai Python (Iris, dendrogram)}

\begin{vidu}[title={Nạp Iris và linkage ward}]
\begin{lstlisting}
import numpy as np
import matplotlib.pyplot as plt
from sklearn.datasets import load_iris
from sklearn.cluster import AgglomerativeClustering
from scipy.cluster.hierarchy import dendrogram, linkage

iris = load_iris()
X = iris.data
y = iris.target

Z = linkage(X, 'ward')
\end{lstlisting}
\end{vidu}

\begin{vidu}[title={Vẽ dendrogram}]
\begin{lstlisting}
plt.figure(figsize=(7.5, 3.5))
plt.title("Iris Dendrogram")
dendrogram(Z)
plt.show()
\end{lstlisting}
\end{vidu}

\tsimg{14_agglomerative/img01.jpg}

\begin{vidu}[title={AgglomerativeClustering và scatter kết quả}]
\begin{lstlisting}
model = AgglomerativeClustering(n_clusters=3)
model.fit(X)
labels = model.labels_

plt.figure(figsize=(7.5, 3.5))
plt.scatter(X[:, 0], X[:, 1], c=labels)
plt.xlabel("Sepal length")
plt.ylabel("Sepal width")
plt.title("Agglomerative Clustering Results")
plt.show()
\end{lstlisting}
\end{vidu}

\tsimg{14_agglomerative/img02.jpg}

\begin{ynghia}[title={ward}]
Phương pháp \texttt{ward} tối thiểu hóa phương sai khoảng cách khi gộp cụm --- thường cho cụm cân bằng hơn trên Iris.
\end{ynghia}

\subsection{Ưu, nhược điểm và ứng dụng}

\begin{tinhchat}[title={Ưu điểm}]
\begin{itemize}
  \item \textbf{Dendrogram} thể hiện quan hệ phân cấp.
  \item Nhiều \textbf{metric} và \textbf{linkage}.
  \item Linh hoạt \textbf{số cụm} khi cắt cây.
  \item Triển khai hiệu quả cho dataset lớn vừa phải.
\end{itemize}
\end{tinhchat}

\begin{luuy}[title={Nhược điểm}]
\begin{itemize}
  \item \textbf{Tốn kém} trên dataset rất lớn.
  \item Cụm \textbf{mất cân bằng} nếu metric/linkage không phù hợp.
  \item Kết quả \textbf{nhạy} lựa chọn metric và linkage.
  \item Dendrogram \textbf{khó đọc} khi có rất nhiều cụm/điểm.
\end{itemize}
\end{luuy}

\begin{ynghia}[title={Ứng dụng}]
\begin{itemize}
  \item \textbf{Phân đoạn ảnh} (Image Segmentation).
  \item \textbf{Phân cụm tài liệu} (Document Clustering).
  \item \textbf{Phân khúc khách hàng} / hành vi khách hàng.
  \item \textbf{Phân khúc thị trường} (Market Segmentation).
  \item \textbf{Phân tích mạng xã hội}.
\end{itemize}
\end{ynghia}
